\documentclass[11pt]{article}
\usepackage{amsmath, amssymb}
\usepackage{geometry}
\geometry{margin=1in}

\title{Routing Notes: Link-State, Distance Vector, and BGP}
\author{}
\date{}

\begin{document}
\maketitle

\section{Overview of Routing}
Routing algorithms compute forwarding tables (destination $\rightarrow$ next hop).

Two main families:
\begin{itemize}
    \item \textbf{Link-State (LS)}: Global topology known
    \item \textbf{Distance-Vector (DV)}: Only neighbor info known
\end{itemize}

\section{Graph Model}
A network is modeled as a graph $G = (V, E)$:
\begin{itemize}
    \item $V$: routers (nodes)
    \item $E$: links (edges)
    \item $c(u,v)$: cost of link
\end{itemize}

\textbf{Goal:} Find shortest path (minimum cost) between nodes.

\section{Link-State Routing (Lecture 15)}
Each router:
\begin{enumerate}
    \item Discovers neighbors
    \item Floods link-state advertisements (LSAs)
    \item Computes shortest paths using Dijkstra
\end{enumerate}

\subsection{Dijkstra's Algorithm}
Maintains:
\begin{itemize}
    \item $D(v)$: current shortest distance from source
    \item $p(v)$: predecessor
\end{itemize}

\textbf{Core idea:} greedily expand shortest known node.

\subsubsection*{Relaxation}
\[
\text{If } D(u) + c(u,v) < D(v), \text{ then update } D(v)
\]

\subsubsection*{Complexity}
\[
O(|V|^2) \quad \text{or} \quad O((|V|+|E|)\log|V|)
\]

\subsection{Shortest Path Tree (SPT)}
Dijkstra builds an SPT rooted at the source.

\textbf{Forwarding table:}
\begin{itemize}
    \item Only first hop matters
    \item Derived from SPT
\end{itemize}

\subsection{OSPF}
Real-world link-state protocol:
\begin{itemize}
    \item Uses LSAs + flooding
    \item Runs directly over IP
    \item Fast convergence (seconds)
\end{itemize}

\subsection{Issues}
\begin{itemize}
    \item High memory (full topology)
    \item Possible oscillations with dynamic costs
\end{itemize}

\section{Distance-Vector Routing (Lecture 16)}
Each router knows:
\begin{itemize}
    \item Cost to neighbors
    \item Neighbors' distance vectors
\end{itemize}

\subsection{Bellman-Ford Equation}
\[
d_x(y) = \min_{v \in \text{neighbors}(x)} \{ c(x,v) + d_v(y) \}
\]

\textbf{Meaning:}
Shortest path = best neighbor + their best path.

\subsection{Distance-Vector Algorithm}
\begin{enumerate}
    \item Initialize distances
    \item Send vector to neighbors
    \item Update using Bellman-Ford
    \item Repeat until convergence
\end{enumerate}

\textbf{Properties:}
\begin{itemize}
    \item Distributed
    \item Iterative
    \item Asynchronous
\end{itemize}

\subsection{Convergence}
\begin{itemize}
    \item Takes at most $|V| - 1$ rounds
    \item Good news propagates fast
    \item Bad news propagates slowly
\end{itemize}

\subsection{Count-to-Infinity Problem}
After link failure:
\begin{itemize}
    \item Nodes form loops
    \item Distances increase gradually
\end{itemize}

\subsection{Fixes}
\begin{itemize}
    \item \textbf{Split Horizon}: don't advertise back to source
    \item \textbf{Poisoned Reverse}: advertise $\infty$
\end{itemize}

\subsection{RIP Protocol}
\begin{itemize}
    \item Metric: hop count
    \item Max: 15 hops
    \item Updates every 30s
\end{itemize}

\section{Link-State vs Distance-Vector}
\begin{tabular}{|c|c|c|}
\hline
 & Link-State & Distance-Vector \\
\hline
Info & Full topology & Neighbor only \\
\hline
Speed & Fast & Slow (failures) \\
\hline
Loops & No & Possible \\
\hline
Example & OSPF & RIP \\
\hline
\end{tabular}

\section{Why One Algorithm Isn't Enough}
The Internet is too large:
\begin{itemize}
    \item Millions of routes
    \item Many organizations
    \item Different policies
\end{itemize}

\section{Autonomous Systems (AS)}
\begin{itemize}
    \item Group of networks under one admin
    \item Each has an ASN
\end{itemize}

\textbf{Types:}
\begin{itemize}
    \item Stub
    \item Multi-homed
    \item Transit
\end{itemize}

\section{BGP (Border Gateway Protocol)}
Inter-AS routing protocol.

\subsection{Key Functions}
\begin{itemize}
    \item Advertise reachable prefixes
    \item Share AS paths
    \item Enforce policies
\end{itemize}

\subsection{Path Vector}
BGP uses:
\begin{itemize}
    \item AS-PATH instead of distance
\end{itemize}

Loop prevention:
\begin{itemize}
    \item Reject routes containing own ASN
\end{itemize}

\subsection{Business Relationships}
\begin{itemize}
    \item Customer $\leftrightarrow$ Provider
    \item Peer-to-peer
\end{itemize}

\textbf{Valley-free routing:}
\[
\text{Customer} \rightarrow \text{Provider} \rightarrow \text{Peer} \rightarrow \text{Customer}
\]

\subsection{Challenges}
\begin{itemize}
    \item Policy over shortest path
    \item Slow convergence (minutes)
    \item Large scale (millions of routes)
    \item Trust-based system
\end{itemize}

\section{Big Picture}
\begin{itemize}
    \item Intra-AS: OSPF (link-state) or RIP (DV)
    \item Inter-AS: BGP (policy-based)
\end{itemize}

\end{document}